\subsection{Simuladores considerados}

Las alternativas evaluadas corresponden a herramientas pertinentes para simulación
cuántica de estado completo, aunque con énfasis distintos en cuanto a rendimiento,
portabilidad, facilidad de uso o experimentación.

\begin{itemize}
    \item \textbf{QuEST (Quantum Exact Simulation Toolkit):} simulador de alto
    rendimiento orientado a vectores de estado y matrices de densidad, con soporte
    para ejecución multihilo, distribuida y acelerada por GPU. Su diseño privilegia
    escalabilidad y desempeño en arquitecturas heterogéneas
    \parencite{Jones2019, quest_repo}.

    \item \textbf{Intel-QS (Intel Quantum Simulator):} simulador de circuitos cuánticos
    basado en representación completa del estado, concebido para aprovechar
    arquitecturas multinúcleo y multinodo. Su enfoque está orientado a entornos HPC y
    a la eficiencia en ejecución paralela \parencite{intelqs_doc}.

    \item \textbf{qsim:} biblioteca de simulación de estado completo integrada al
    ecosistema de Cirq, optimizada para calcular el estado final de circuitos cuánticos.
    Su organización del código distingue componentes para simulación y manejo del
    vector de estado, con énfasis en ejecución eficiente
    \parencite{qsim_repo, qsim_cirq}.

    \item \textbf{Quantum++:} biblioteca general en C++ para simulación cuántica,
    distribuida como librería \emph{header-only} y diseñada con énfasis en portabilidad,
    facilidad de uso y rendimiento. Su arquitectura resulta adecuada para prototipado y
    experimentación ligera \parencite{qpp_paper}.

    \item \textbf{TMFQSfullstate:} prototipo de simulador desarrollado en C++ con foco
    en estudiar estrategias de representación y consumo de memoria sobre simulación
    de estado completo. Su diseño favorece la intervención directa sobre la
    representación del estado y la exploración de variantes relacionadas con gestión de
    memoria y paralelismo
    \parencite{diaz2024software, diaz_mdpi_mmss, tmfqs_repo}.
\end{itemize}

En conjunto, las alternativas cubren dos perfiles complementarios. QuEST, Intel-QS y
qsim representan herramientas con fuerte orientación a rendimiento y ejecución
eficiente. Quantum++ y TMFQSfullstate, en cambio, ofrecen condiciones más
favorables para intervención directa sobre el código y experimentación sobre la
representación del estado.
